\chapter{Requirements Analysis and System Design}

\section{Introduction}
This chapter formalizes the functional and non-functional requirements of \projectname{} and translates them into an architectural design.

\section{Functional Requirements}
The platform must support heterogeneous data ingestion, distributed storage, distributed transformation, catalog integration, SQL access, orchestration, monitoring, and validation workflows.

\section{Non-Functional Requirements}
The platform must demonstrate scalability, modularity, security, observability, service separation, reproducibility, and suitability for academic demonstration.

\section{Actors and Usage Perspective}
This section should identify key actors such as data engineers, analysts, platform administrators, and future data science consumers.

\section{Global Logical Architecture}
This section should present the logical data flow from ingestion to validation and analytical exposure.

\section{Physical Multi-Node Architecture}
This section should present the role separation across the control plane, VM1, VM2, and VM3.

\section{Design Rationale}
The final design should be explained as a balance between academic rigor, distributed execution proof, security hardening, and resource constraints.

\section{Chapter Conclusion}
This chapter provides the design baseline for the concrete implementation detailed in the next chapter.

